\section{Pemodelan Dependensi}

Kesenjangan antara efisiensi memori dan kemampuan memodelkan dependensi mendorong penggunaan model probabilistik yang tetap mampu menangkap hubungan antar-variabel.
\textit{Bayesian Optimization Algorithm} (BOA) merupakan EDA multivariat yang mempelajari \textit{Bayesian network} tanpa membatasi derajat dependensi antar-variabel \parencite{pelikan_iboa_2008}.
Strukturnya dipelajari secara \textit{greedy} dari graf kosong dengan menambahkan \textit{edge} yang paling meningkatkan skor \textit{Bayesin Information Criterion} (BIC)\label{first:bic} hingga tidak diperoleh peningkatan lebih lanjut, sedangkan parameter jaringan diestimasi dari populasi solusi terpilih.
BOA tetap menjadi salah satu pendekatan EDA multivariat yang digunakan dengan berbagai bentuk model, termasuk jaringan Bayes diskret, Gaussian, semiparametrik, dan berbasis \textit{kernel density} \parencite{soloviev_edaspy_2024}.

Solusi baru pada BOA dibangkitkan melalui \textit{logic sampling} berdasarkan struktur jaringan yang telah dipelajari \parencite{pelikan_iboa_2008}.
Variabel terlebih dahulu diurutkan secara \textit{ancestral} sehingga setiap induk mendahului variabel anaknya, kemudian nilainya di-\textit{sampling} berdasarkan distribusi bersyarat terhadap nilai induk yang telah ditentukan.
Mekanisme ini memungkinkan kebutuhan komputasi dan memori seiring kompleksitas jaringan.

Di antara asumsi independensi penuh EDA univariat dan dependensi multivariat pada BOA, terdapat pendekatan bivariat berbasis \textit{dependency tree} \textcite{chow_approximating_1968}.
Metode ini mengaproksimasi distribusi gabungan multidimensi melalui struktur pohon yang hanya mempertahankan dependensi berpasangan.
Struktur tersebut diperoleh dengan menghitung \textit{mutual information} setiap pasangan variabel dan membentuk \textit{maximum spanning tree}.
\textcite{chow_approximating_1968} menunjukkan bahwa pohon yang dihasilkan memberikan aproksimasi dengan divergensi Kullback-Leibler minimum terhadap distribusi sebenarnya.

\begin{figure}[H]
    \centering
    \includegraphics[width=0.9\textwidth]{gambar_dependensi.png}
    \caption{Spektrum Pemodelan Dependensi pada EDA dan Posisi CoDA-EDA}
    \label{fig:spektrum-dependensi}
\end{figure}

Gambar \ref{fig:spektrum-dependensi} menunjukkan adanya \textit{trade-off} antara kemampuan memodelkan dependensi dan biaya memori/komputasi.
CoDA-EDA ditempatkan di antara pendekatan univariat dan multivariat penuh dengan mempertahankan dependensi antar-atribut yang relevan, tetapi dengan struktur yang dibatasi agar kebutuhan sumber dayanya tetap rendah.

Seperti EDA kanonik, BOA standar masih membutuhkan populasi eksplisit untuk mempelajari struktur dan parameter \textit{Bayesian network}.
\textcite{pelikan_iboa_2008} mengembangkan \textit{incremental BOA} (iBOA) yang menghilangkan kebutuhan tersebut dengan memperbarui struktur dan parameter jaringan secara inkremental melalui skema turnamen pemenang/pecundang.
Dengan demikian, iBOA menunjukkan bahwa pemodelan dependensi multivariat dapat dilakukan tanpa menyimpan populasi secara eksplisit.

Namun, skalabilitas iBOA tetap menjadi tantangan.
Penyimpanan tabel probabilitas marginal untuk berbagai kemungkinan struktur dependensi dapat membutuhkan hingga $O(n^k)$\label{first:symbol-onk} tabel, dengan $k$\label{first:symbol-k} sebagai orde dependensi maksimum yang diizinkan \parencite{pelikan_iboa_2008}.
Tanpa pembatasan struktur, kebutuhan ini dapat mengurangi keunggulan memori pendekatan \textit{compact}.
Celan tersebut menjadi dasar rancangan CoDA-EDA, yang menggabungkan pembaruan inkremental tanpa populasi eksplisit dengan struktur dependensi yang dibatasi agar tetap efisien secara memori dan komputasi.